\section{Verification-axis identifiability and attainability}
\label{sec:axis-theory}

A benchmark does not measure a scientific competence merely because it assigns a
label and reports an accuracy.  The label must be recoverable from the intended
scientific information, unavailable from construction nuisances, expressible by
the systems being compared, and capable of varying over the comparison panel.
These are different requirements.  We formalize them before interpreting the
protected results.

Let $Y_a\in\mathcal{T}$ be the correct terminal on verification axis $a$ for a
random benchmark unit, let $D$ be the representation available to a decision
rule, and let finite nonempty $\mathcal{A}\subseteq\mathcal{T}$ be the terminals
that the rule is allowed to emit.  For a finite benchmark distribution, the
smallest attainable exact-terminal error is
\begin{equation}
  R_a^*(D,\mathcal{A})
  = 1-\mathbb{E}_{D}\!\left[
      \max_{t\in\mathcal{A}}\Pr(Y_a=t\mid D)
    \right].
  \label{eq:bayes-axis-risk}
\end{equation}
This quantity separates a poor implementation from a task that its information
or interface makes impossible.

\paragraph{Proposition 1 (exact axis attainability).}
There exists a decision rule $g(D)\in\mathcal{A}$ with zero exact-terminal error
if and only if every positive-probability fibre of $D$ contains one target
terminal and that terminal belongs to $\mathcal{A}$.  Equivalently,
$Y_a=g(D)$ almost surely for some $g$ with codomain $\mathcal{A}$.  This is an
almost-sure benchmark statement: it says nothing pointwise about worlds on null
fibres.

\emph{Proof.}  On a fibre $D=d$, any rule must emit one terminal.  It is correct
on the whole fibre exactly when the conditional support of $Y_a$ is the
singleton containing that terminal.  Summing the best fibrewise choices gives
Equation~\ref{eq:bayes-axis-risk}, and zero risk gives the stated condition.
\hfill$\square$

\paragraph{Corollary 1 (donor-product factorization at two authorities).}
Let $D$ collect the states exposed by a product of provenance, verification,
custody, freshness, and generic-authorization donors, and let $Y_{\rm promote}$
be the target scientific-promotion terminal. Pointwise extensional equivalence
on a declared world class holds if and only if the target terminal is constant
on every donor-state fibre in that class and belongs to the output alphabet.
Zero error under a benchmark distribution requires only the corresponding
condition on positive-probability fibres; null fibres cannot be certified by a
distributional score. If a positive-probability fibre mixes target terminals,
every donor-only rule has error at least
\begin{equation}
  \sum_d \Pr(D=d)
  \left[1-\max_{t\in\mathcal{A}}
  \Pr(Y_{\rm promote}=t\mid D=d)\right].
  \label{eq:mixed-fibre-bound}
\end{equation}
The pointwise statement is the set-theoretic factorization argument; the
distributional statement is Proposition~1, and the displayed lower bound is
attained by a fibrewise Bayes rule on the finite benchmark. Thus a separate module has no inherent advantage once a donor product receives
a target-sufficient representation and the same decision relation; conversely,
no optimization of a representation that merges differently licensed states
can recover the missing distinction.  This is a factorization boundary, not a
claim that scientific authority must be centralized.

\paragraph{Corollary 2 (terminal-preserving comparator adapters).}
Let $V\in\mathcal V$ be the complete candidate-visible record supplied to an
external comparator and let $X\in\mathcal X$ be its native output, where
$\mathcal V$ and $\mathcal X$ are standard-Borel spaces.  Allow $X$ to be
generated by a measurable randomized kernel from $V$, so that the comparator
output adds no world information beyond $V$.  A benchmark adapter is a
prospectively declared measurable map $h:\mathcal X\to\mathcal A$; it is not a
post-hoc relabelling rule.  Its exact-terminal
risk is bounded below by
\[
  R_a^*(X,\mathcal A)=1-\mathbb E_X\!\left[
    \max_{t\in\mathcal A}\Pr(Y_a=t\mid X)
  \right]
  \ \ge\ R_a^*(V,\mathcal A).
\]
The first bound is attained by a measurable fixed-order Bayes adapter, and the
second inequality is the data-processing consequence of $Y_a-V-X$.  Thus a
native output is admissible for a three-way exact-axis endpoint with zero
benchmark error if and only if its positive-probability output fibres are
target-pure and their targets belong to $\mathcal A$.  For a deterministic
native-output map, pointwise admissibility on a declared world class requires
the same purity on every native-output fibre, including benchmark-null fibres.
In particular, a binary native
alphabet cannot attain an endpoint on which all three target terminals have
positive probability, and a native \textsc{Block}, ``not attributable,'' parse
failure, or free-text insufficiency response cannot be renamed
\textsc{CannotCheck} unless the predeclared semantics-preserving map satisfies
this fibre condition.  When it does not, the mixed-fibre risk is an
interface-attainability result and must not be reported as comparator
inferiority.

\paragraph{Proposition 2 (claim-level indistinguishability).}
Let $W_0$ and $W_1$ be two scientific worlds in which a proposed competence
claim has opposite truth values, and let $O$ be the complete observable record
available to an adjudicator.  Write $P_i=\mathcal L(O\mid W_i)$ for probability
measures on one common measurable record space $(\mathcal O,\mathcal F_O)$, and
use the convention
$\operatorname{TV}(P_0,P_1)=\sup_{B\in\mathcal F_O}|P_0(B)-P_1(B)|$.
Under equal prior weight, every measurable, possibly randomized adjudicator has
error at least
\begin{equation}
  \frac{1-\operatorname{TV}
  (\mathcal{L}(O\mid W_0),\mathcal{L}(O\mid W_1))}{2}.
  \label{eq:claim-identifiability-bound}
\end{equation}
The infimum error over adjudicators equals this bound.  For two probability
measures it is attained: a Hahn decomposition of the finite signed measure
$P_1-P_0$ supplies a maximizing acceptance region.  In particular, if the
observable records have the same distribution, the competence claim is not
identifiable and the error is exactly $1/2$ for every adjudicator.

\emph{Proof.}  For any acceptance region $B$, the equal-prior error is
$\{P_0(B)+P_1(B^{\mathsf c})\}/2$.  Taking the infimum over $B$ gives
one half of one minus the total-variation distance.  The positive set in a
Hahn decomposition of $P_1-P_0$ attains the supremum, while randomized
acceptance probabilities cannot improve on it because their risks are mixtures
of deterministic risks. \hfill$\square$

This second bound concerns interpretation rather than terminal prediction.  A
score can separate output interfaces while leaving two explanations of that
score observationally equivalent. For example, one rule may recognize an
epistemic insufficiency, whereas another may be forced into the only available
non-promoting terminal.  A broader competence headline remains
unidentified until a discriminator changes the observable law of those worlds.

\paragraph{Nuisance identifiability.}
Let $N$ contain construction-only quantities such as list length, missingness,
template, key order, path name, or label-bearing metadata.  The nuisance-only
Bayes advantage is
\begin{equation}
  I_a(N)=\mathbb{E}_{N}\!\left[\max_t\Pr(Y_a=t\mid N)\right]
  -\max_t\Pr(Y_a=t).
\end{equation}
A scientific interpretation requires the intended semantic representation to
make the terminal attainable while nuisance information contributes no material
advantage.  A finite shortcut register cannot certify $I_a(N)=0$ for every
possible decoder.  It certifies only that the frozen probe class found no such
advantage, which is why claim-axis authorization must name the probe class and
cannot be promoted to whole-register or naturalistic clearance.

\paragraph{Panel resolution and interface attainability.}
For a frozen panel $\mathcal{H}$ and score $M_a$, define its observed resolution
as
\[
  \rho_a(\mathcal{H})=\max_{h\in\mathcal{H}}M_a(h)
  -\min_{h\in\mathcal{H}}M_a(h).
\]
When $\rho_a=0$, the experiment supplies no comparative ordering on that axis;
it does not establish equality of the underlying competences.  When
$\rho_a>0$ but some comparators lack a target terminal in their output alphabet,
the contrast establishes interface expressiveness before it establishes the
finer judgement of when that terminal is scientifically appropriate.  A broad
verification-axis claim is therefore licensed only inside the intersection of
four conditions: semantic attainability, nuisance resistance, matched terminal
attainability, and nonzero panel resolution.  Failures fall into distinct
research programmes rather than one generic null: repair the construction,
enlarge the interface, expose a target-sufficient representation, or improve
the decision rule.

\paragraph{Application to the three P4 evidence layers.}
The historical V2 H3 axis had $\rho_a=0$: an empty evidence list exposed the
gold family and every panel member scored 30/30.  That result remains
\textsc{NotSupported}.  The later V3 construction removes the registered
shortcuts and obtains nonzero resolution, but nine of ten comparator mechanisms
cannot emit \textsc{CannotCheck}; its exact claim is therefore terminal
expressiveness under the frozen gate lattice, not general epistemic judgement.

Stated as an estimand: where a comparator interface cannot emit
\textsc{CannotCheck}, the quantity this comparison estimates is \emph{terminal
attainability} on the axis, and nothing else. A system that cannot express the
correct terminal cannot be observed to reach it, so a gap measured against such
an interface is a fact about what the interface can say, not about what the
system knows. Reading it as epistemic competence would attribute to judgement
what belongs to expressiveness.
The P4-X exact-contract successor tests the donor-factorization boundary more
directly.  Its ideal typed donor product receives the target-sufficient
coordinates and predicate and ties exactly at 400/400, while the registered
generic-authorization and compensatory products do not.  These layers diagnose
different limits and are not pooled.
